% Part 2 chapter carrying "Lensing-Dynamics Mass Discrepancy as a Redshift-Dependent Signature of Informational Gravity" (Mahaffey, 25 Feb 2026)
% in its order. Read in full 2026-10-02 (PDF text 307 lines in this extraction, 316 in the check file's; every line read).
% Numbers: docs/verification/scripts/verify_obs_chapters.py (sections A, E) and verify_s8_trend.py (item 5).
% Check: docs/verification/observations/LENSING_DYNAMICS_CHECK.md; errata LD1-LD5. Carried in the Level 1 form only; the test is on hold.
\chapter{Lensing mass and dynamical mass}\label{ch:lensdyn}

\section*{The place}
A galaxy cluster: about $10^{14}$--$10^{15}\,M_\odot$ of matter bound on timelike worldlines, crossed by light from galaxies behind it. Matter in the
cluster responds to the Newtonian potential $\Psi$; light responds to the lensing potential $(\Phi+\Psi)/2$. If the informational term changes the first
and not the second, the two mass estimates of one cluster differ, and the difference depends on when the light left. Whether the term does change
$\Psi$ inside a cluster depends on its implementation form, which is not yet settled (Section~\ref{sec:ld_hold}).

\section{The ratio in the Level 1 form}
In the Level~1 form the informational term enters the Poisson equation for $\Psi$ as an effective Newton constant $\mu G$
(Chapter~\ref{ch:latetime}):
\begin{equation}
\nabla^2\Psi=-4\pi G a^2\mu(a)\bar\rho\delta,\qquad\nabla^2(\Phi+\Psi)=-8\pi Ga^2\Sigma(a)\bar\rho\delta,\qquad\Sigma=1.
\end{equation}
A dynamical mass (galaxy velocities or gas in hydrostatic equilibrium) is read from $\Psi$, so $M_{\rm dyn}=\mu M_{\rm true}$; a lensing mass is read from
$\Phi+\Psi$, so $M_{\rm lens}=M_{\rm true}$. Then
\begin{equation}
\frac{M_{\rm lens}}{M_{\rm dyn}}=\frac{\Sigma}{\mu}=\frac{1}{\mu(z)}=1+\frac{\beta_m\,e^{1-1/a}}{\Omega_ma^{-3}+\Omega_\Lambda},
\label{eq:ld_ratio}\end{equation}
with $\beta_m=\Omega_m/2$ and nothing adjusted. \calc
\begin{center}\small\begin{tabular}{lcccccccccc}\toprule
$z$ & 0 & 0.1 & 0.2 & 0.3 & 0.5 & 0.7 & 1.0 & 1.5 & 2.0 & 3.0\\\midrule
$\mu$ & 0.864 & 0.886 & 0.905 & 0.922 & 0.948 & 0.966 & 0.982 & 0.994 & 0.998 & 1.000\\
$M_{\rm lens}/M_{\rm dyn}$ & 1.158 & 1.129 & 1.105 & 1.085 & 1.055 & 1.035 & 1.018 & 1.006 & 1.002 & 1.000\\\bottomrule
\end{tabular}\end{center}
The ratio is $15.8\,\%$ above unity today and falls monotonically with redshift, with slope $dR/dz=-0.31$ at $z=0$ and $-0.18$ at $z=0.3$. It
depends on the cluster's redshift only, not on its mass, radius or dynamical state. \calc

\section{The form of the term}\label{sec:ld_hold}
\authorhold{The ratio \eqref{eq:ld_ratio} needs $\mu$ in the Poisson equation for $\Psi$ (the Level~1 effective-coupling form). In the Level~2 form the
term is a matter-sector expansion rate in the perturbation equations, and the linearised Einstein equations for the potentials are unmodified
(Chapter~\ref{ch:level2}). Then $\Psi$ obeys the standard Poisson equation, matter in a virialised cluster moves in the true potential, and
$M_{\rm lens}/M_{\rm dyn}=1$ at every redshift. The source paper states both forms. The prediction is carried here in the Level~1 form only, and only
until two questions are settled: which form is the physical one, and whether a linear $\mu$ applies inside virialised, non-linear clusters at all.}
\openprob

\section{Other explanations and how they differ}
\begin{itemize}
\item \emph{Hydrostatic bias.} Non-thermal pressure (turbulence, bulk motion) makes X-ray masses low. Simulations find the non-thermal fraction rising
with redshift, as younger clusters are less relaxed~\cite{Nelson2014,ShiKomatsu2014}; the bias then grows with $z$, the opposite slope to
Eq.~\eqref{eq:ld_ratio}. \observed
\item \emph{MOND-type theories.} The discrepancy depends on the internal acceleration profile of each cluster, not on its redshift.
\item \emph{$f(R)$ gravity.} In the quasi-static limit $f(R)$ has $\Sigma=1$ and $1\le\mu\le4/3$~\cite{PogosianSilvestri2016}: the lensing mass equals
the true mass and the dynamical mass is larger, $M_{\rm lens}/M_{\rm dyn}\le1$. The self-accelerating branch of DGP shares the signs $\mu<1$, $\Sigma=1$,
but carries a ghost and is excluded. The signs are therefore not unique to IAM; its content is the parameter-free curve \eqref{eq:ld_ratio}.
\end{itemize}

\section{What the data say now}
\begin{itemize}
\item \emph{Planck SZ cluster counts.} With the primary-CMB cosmology, the counts need a mass bias $1-b=0.58\pm0.04$, against the baseline
$1-b=0.80$ suggested by simulations~\cite{Planck2015SZ}. \observed The informational term lowers $\sigma_8$ by $1.5\,\%$ (Level~1 chains, $0.8139\to0.8014$) or $1.1\,\%$ (Level~2,
$0.8087\to0.7998$). That reduces the tension; it does not remove it. \calc
\item \emph{Lensing calibrations of cluster masses.} The Canadian Cluster Comparison Project~\cite{Hoekstra2015} and Weighing the
Giants~\cite{vonderLinden2014} calibrate Planck and X-ray masses against weak lensing. Each publishes its result in its own convention ($1-b$, or
$M_{\rm Planck}/M_{\rm WL}$, within its own aperture), and X-ray hydrostatic masses also follow $\Psi$. A comparison with Eq.~\eqref{eq:ld_ratio} must
take each ratio from its source table and match its redshift distribution; this book does not yet print one as a test. \openprob
\end{itemize}

\section{The test}
\prediction\ In the Level~1 form, $M_{\rm lens}/M_{\rm dyn}$ in at least five redshift bins over $0.1<z<2$ with one mass method throughout (weak lensing
against spectroscopic velocity dispersions, the same profile model for both) follows Eq.~\eqref{eq:ld_ratio}: $1.105$ at $z=0.2$, $1.055$ at $0.5$,
$1.018$ at $1$, $1.002$ at $2$. Euclid, Rubin and eROSITA provide the samples. A fit compares three shapes: a constant (hydrostatic bias), Eq.~\eqref{eq:ld_ratio}
with no free parameter, and a two-parameter power law $1+A(1+z)^n$. The discriminant is the shape in redshift, not the existence of an offset.
\begin{itemize}
\item Falsified (in this form) by a ratio flat in $z$, by one rising with $z$, or by a $z=0$ value away from 1.158 at more than $3\sigma$ in a clean
sample.
\item Lensing against velocity dispersions and lensing against X-ray masses are the same test only in the Level~1 form.
\end{itemize}

\section{Corrections applied}
The prediction is carried with its form stated and on hold. $f(R)$ has $\Sigma=1$ in the quasi-static limit, and self-accelerating DGP shares the signs
of $\mu$ and $\Sigma$; the uniqueness claim is replaced by the parameter-free shape. The SZ tension is reduced, not resolved. The published cluster
ratios are not printed until traced to their tables. Chain counts and $\sigma_8$ are from the final chains. Appendix~\ref{app:errata} lists each item.

\section{Status}
\begin{center}\small\begin{tabularx}{\linewidth}{>{\raggedright\arraybackslash}Xl}\toprule
Result & Label\\\midrule
$M_{\rm lens}/M_{\rm dyn}=1/\mu(z)$, Eq.~\eqref{eq:ld_ratio}, Level~1 form & \calc\ \prediction\\
Which implementation form applies inside clusters & \openprob\\
Non-thermal pressure slope positive in $z$ & \observed\\
$\sigma_8$ shift and the SZ counts & \calc\\
Comparison with published cluster ratios & \openprob\\\bottomrule
\end{tabularx}\end{center}
